\documentclass[11pt,a4paper]{article}

\usepackage{fontspec}
\usepackage{polyglossia}
\setmainlanguage{russian}
\setotherlanguage{english}
\setmainfont{PT Serif}[Ligatures=TeX]
\setsansfont{PT Sans}[Ligatures=TeX]
\setmonofont{PT Mono}[Scale=0.92]

% macOS ставит ParaType как .ttc, и fontconfig не рапортует покрытие кириллицы;
% polyglossia верит fontconfig и отказывается принимать основной шрифт, хотя
% кириллица в нём есть (проверено отдельной сборкой). Объявляем семейства явно ---
% это штатный обход, рекомендованный самим polyglossia, и на Linux он безвреден.
\newfontfamily\cyrillicfont{PT Serif}[Ligatures=TeX]
\newfontfamily\cyrillicfontsf{PT Sans}[Ligatures=TeX]
\newfontfamily\cyrillicfonttt{PT Mono}[Scale=0.92]

\usepackage{amsmath,amssymb}
\usepackage{graphicx}
\usepackage{booktabs}
\usepackage{tabularx}
\usepackage[table]{xcolor}
\usepackage[margin=27mm,top=25mm,bottom=27mm]{geometry}
\usepackage{caption}
\usepackage{enumitem}
\usepackage{hyperref}
\usepackage{microtype}

\captionsetup{font=small,labelfont=bf,labelsep=period,justification=justified}
\captionsetup[table]{skip=4pt}
\setlist{itemsep=2pt,parsep=0pt,topsep=4pt}
\hypersetup{colorlinks=true,linkcolor=black,urlcolor=black,citecolor=black,
            pdftitle={Модель предсказания ингибирования цитохромов P450},
            pdfauthor={команда CYP}}

\emergencystretch=2.2em
\setlength{\parskip}{0pt}
\setlength{\parindent}{1.2em}
\linespread{1.06}
\renewcommand{\arraystretch}{1.12}

\newcommand{\pic}{\mathrm{pIC}_{50}}
\newcommand{\pC}{\mathrm{p}C}
\newcommand{\Var}{\mathrm{Var}}
\newcommand{\E}{\mathbb{E}}

\title{\vspace{-14mm}\textbf{\Large Модель предсказания ингибирования цитохромов P450}\\[3mm]
\normalsize Устройство решения для OpenADMET CYP Blind Challenge: что в данных,\\
что мы измерили, как устроена модель и почему именно так}
\author{}
\date{\normalsize 27 августа 2026}
